
\chapter{Introduction}
\label{ch:1}

\section{Background}
\label{sec:1.1}
Detectors of abusive language are deployed under latency and memory budgets that exclude the large
encoders that score best on the benchmarks \cite{schwartz2020green, patterson2021carbon,
zhou2019edge}. Knowledge distillation is the standard response \cite{hinton2015distilling,
bucilua2006model}: a compact student is trained to imitate a large teacher's softened outputs, and
in the multi-teacher variants this literature increasingly favours, the outputs of several teachers
combined by a weighting that trusts each where it is reliable \cite{you2017learning, wu2021one,
zhang2022confidence, yuan2021reinforced}. In abusive-language detection the recipe is almost always
applied in one way: the teachers are fine-tuned on the task's labelled training split, and the
student is distilled on that same split \cite{wu2021one, prasomphan2025mtkd}. The promise is that a
committee of teachers with complementary expertise holds more of what the task requires than any one
of them, and that the student, small enough to run on a phone, receives it.

This thesis continues the Pre-Thesis II report of February 2026, which proposed Dynamic Multi-Teacher
Homogeneous Knowledge Distillation (D-MTHD) and evaluated it on the Wikipedia Talk personal-attacks corpus
with two teachers and a DistilBERT student. That report found the student at 0.8947 macro-F1 against
teachers at 0.8936 and 0.8904, from one run of each configuration, and named two tasks for this phase:
extending the method to sarcastic and implicit cyberbullying, and testing it against a student trained
without teachers. This thesis does both, on a cleaned fine-grained cyberbullying corpus, and what it
finds changes the method.

\section{Motivation}
\label{sec:1.2}
Moderation at the scale of a social platform depends on detectors cheap enough to run on every post,
and increasingly on the device where the post is written \cite{zhou2019edge, schwartz2020green}.
Compact models make that possible, and distillation is how compact models are usually made
competitive. Two facts make the question worth a thesis. The first is that the abuse detectors miss
most is the abuse that is implied rather than stated: every survey of cyberbullying detection names
sarcastic and indirect abuse as the open case \cite{rosa2019automatic, salawu2020approaches,
emmery2021current}, and implicit hate has needed benchmarks of its own \cite{elsherief2021latent,
ocampo2023indepth}. The second is that the distillation methods proposed for this task, and the
Pre-Thesis II version of this one, have been reported without the controls that would show whether
the teachers help at all. A practitioner choosing how to train a compact detector needs to know
whether distillation moves it, where the gain comes from, and whether any of it reaches implied
abuse. This thesis sets out to answer those three questions with measurements rather than
assumptions.

\section{Problem Statement}
\label{sec:1.3}
Three problems meet in this task. The first is deployment: the encoders that detect abuse best are too
large and slow for on-device or high-throughput moderation, and compact students trail them. The second
is implication: the abuse that detectors and annotators miss most is carried by irony, stereotype and
coded reference, and the widely used benchmarks do not label it. The third is evidence: the methods
proposed to close the first gap, committees of teachers weighted per instance, are usually reported
without the control that would show whether the teachers helped at all, which is the same student
trained without a teacher, and without repeated runs or confidence intervals. The Pre-Thesis II report
shared this gap. A distillation method that is not shown to beat that control cannot be said to work, and
a gain that is not measured on implicit abuse cannot be said to address it.

\section{Research Questions and Objectives}
\label{sec:1.4}
The thesis asks four questions.
\begin{description}
  \item[RQ1.] Does multi-teacher distillation with per-instance reliability weighting improve a compact
    cyberbullying detector over the same student fine-tuned without teachers?
  \item[RQ2.] If it does not, why not?
  \item[RQ3.] Does distillation on unlabelled text the teachers have never seen improve the student, and
    how does the gain depend on the amount and composition of that text, on the number of optimisation
    updates, and on the student?
  \item[RQ4.] Does any of these interventions improve the student's ability to tell implied abuse from
    harmless sarcasm?
\end{description}
To answer them, the thesis sets four objectives.
\begin{enumerate}
  \item Build a clean benchmark and a controlled grid of five students, three committees, four objectives
    and three seeds, with a no-teacher control and paired bootstrap intervals for every comparison.
  \item Measure the mechanism: what the teachers have memorised, how far they agree, and how much any
    weighting of them could gain.
  \item Construct an unlabelled transfer set screened against every evaluation set, and measure the
    out-of-sample gain against its size and composition and against fine-tuning at matched compute,
    with the predictions written before each run.
  \item Measure the discrimination of implied abuse without a threshold, on probe sets and on a benchmark
    that labels implication.
\end{enumerate}

\section{Approach and Summary of Findings}
\label{sec:1.5}
We tested that promise under the controls it is rarely given, and then tested the alternative the
controls pointed to. The first half of the paper is an audit. We built the method the literature
suggests: a committee of task-adapted specialist teachers, a general encoder, an abusive-language
specialist and an irony specialist, weighted per instance by reliability, with an auxiliary irony
head and a fourth teacher trained on a corpus in which implication is a label. We measured it on a
cleaned benchmark against the baseline these papers rarely include, the same student trained with no
teacher at all, over five students, three committees, every weighting temperature from 0.05 to 5 and
three seeds, with every difference tested by paired bootstrap. Nothing moved. Distillation raised
every student by 0.001 to 0.007 macro-F1 and no interval excluded zero; per-instance weighting never
differed from uniform averaging; the specialist added nothing in a committee, alone or as
pre-training. And the reason is measurable rather than argued: on the split the teachers were
fine-tuned on, their training losses end at 0.03 to 0.08, so their soft labels are the gold labels,
the reliability signal that weights them is saturated, and distillation on that split reduces to
fine-tuning with a smoothed target.

The second half follows from the first. If soft labels carry information only where the teacher is
uncertain, the student should be distilled on text the teacher has never fitted. We added an
unlabelled transfer set of tweets, screened against every split, probe and benchmark, on which the
student trains from the teacher's soft labels alone, and measured its effect the way the audit had
taught us to: with the committee against a single teacher, with a corrected out-of-sample
reliability weighting against averaging, with hard pseudo-labels against soft, with the size of the
set varied over six nested subsets from 5,000 to 168,000 tweets, with generic tweets against
abuse-related ones at fixed size, with fine-tuning alone run for the same number of updates, and
with the largest set on two more students. The predictions were written before each run and scored
by a script against the tables. The gain exists, grows monotonically with the amount of text,
reaches +0.012 macro-F1 on the 11M student with an interval clear of zero, comes from generic tweets
nearly as well as from abuse-related ones, survives fine-tuning at matched compute, which gains
nothing on its own, and holds on a 10M BiLSTM (+0.016) and a 29M BERT (+0.006). One teacher does it
as well as three; the weighting, the committee and the specialist add nothing here either.
Distilling a compact student on an unlabelled transfer set is not new \cite{hinton2015distilling,
turc2019wellread, tang2019distilling}; what is new is the controlled demonstration, in this
application, that it is the only lever that moves the student, and the measurement of what it moves.

The knowledge that matters most in this task is also the one that does not move. Abuse that names
its target is largely solved by lexical models; abuse carried by implication, stereotype, irony or
coded reference is where detectors and annotators both fail \cite{elsherief2021latent,
ocampo2023indepth, hartvigsen2022toxigen}, every survey of cyberbullying detection names it as the
open case \cite{rosa2019automatic, salawu2020approaches, emmery2021current}, and no widely used
benchmark labels it \cite{wang2020sosnet, wulczyn2017exmachina}. We measured it with a
threshold-free metric on held-out probes of ironic abuse and harmless sarcasm, and with a corpus
that labels implication. Across 184 models in the grid, recall on ironic abuse and the
false-positive rate on harmless sarcasm rise together; the ability to tell them apart is set by
pre-training and unmoved by every objective, committee, teacher and transfer set we tried. The
out-of-sample gain is a task gain. What the student learns from the unlabelled text is the boundary
between the two classes the labelled split draws worst, not the reading of implication.

\section{Contributions}
\label{sec:1.6}
The contributions are as follows.

\begin{enumerate}
  \item \textbf{A controlled audit of multi-teacher distillation for abusive language}
    (Sections~\ref{sec:5.2} to~\ref{sec:5.6}). Five students across three architecture families,
    three committees, four objectives, a specialist teacher trained on labelled implication,
    weighting temperatures over a hundredfold range, three seeds, 100 paired comparisons with
    bootstrap intervals. In sample, nothing beats fine-tuning by more than the test set can resolve,
    and the mechanism is measured: memorised labels, saturated reliability, and diversity spent by
    the task adaptation that makes teachers comparable.
  \item \textbf{The out-of-sample result and its curve} (Sections~\ref{sec:5.7} and~\ref{sec:5.8}).
    Soft labels on unlabelled tweets the teachers have not fitted raise the student, monotonically
    over six sizes to +0.012 at 168,000 rows, from any tweets of the platform, with fine-tuning at
    matched compute gaining nothing, on three compact students of two families; the committee, its
    weighting and hard pseudo-labels add nothing.
  \item \textbf{A threshold-free protocol for implication} (Sections~\ref{sec:4.6}
    and~\ref{sec:5.9}), which shows that recall at a fixed cut measures readiness to fire and that
    discrimination of implied abuse is fixed by pre-training, on the probes and on a corpus that
    labels implication.
  \item \textbf{A single-source implicit-abuse benchmark and a screened transfer-set construction}
    (Sections~\ref{sec:4.2} and~\ref{sec:4.4}), with the evidence that motivates each: pooling the
    two public implicit-hate corpora yields a benchmark solvable by recognising the source, and a
    transfer set must be screened against every evaluation set, including the held-out corpus one of
    its sources is incorporated in.
\end{enumerate}

We do not claim a new method. The weighting we audited is the one this literature keeps proposing,
the controls we ran are the ones it keeps omitting, and the recipe we recommend is the oldest one in
distillation, measured for the first time in this domain against the alternatives and with the
confounds removed.

\section{Thesis Organisation}
\label{sec:1.7}
Chapter~\ref{ch:2} reviews the literature on cyberbullying and implicit abuse, on distilling compact
language models, on multi-teacher distillation, and on evaluation practice. Chapter~\ref{ch:3} describes
the audited method and the out-of-sample setting. Chapter~\ref{ch:4} describes the data, the probe sets,
the transfer set and the protocol. Chapter~\ref{ch:5} reports the results in the order of the argument.
Chapter~\ref{ch:6} answers the research questions, discusses what the findings mean and records how the
study changed. Chapter~\ref{ch:7} states the limitations and future work, and Chapter~\ref{ch:8}
concludes.
